\section*{الفصل \ref{ch:g2:mixing-and-dissolving} --- الخلط والذوبان}

\begin{solution}{exo:g2:mixing-and-dissolving:1}
السكر والملح يذوبان. أما الرمل والطباشير والحصى فلا.
\end{solution}

\begin{solution}{exo:g2:mixing-and-dissolving:2}
\omterm{def:g2:mixing-and-dissolving:dissolve}{المحلول} هو السائل الرائق الذي نحصل عليه حين \omterm{def:g2:mixing-and-dissolving:dissolve}{يذوب} جسم صلب في الماء.
مثالان: الماء المحلّى، والماء المالح (والشاي المحلّى \omterm{def:g2:mixing-and-dissolving:dissolve}{محلول} أيضًا).
\end{solution}

\begin{solution}{exo:g2:mixing-and-dissolving:3}
نعم: وُضع الشوفان والزبيب والجوز والبذور معًا. ولا يزال كل جزء يُرى
ويمكن التقاطه.
\end{solution}

\begin{solution}{exo:g2:mixing-and-dissolving:4}
لا. الجسم الصلب الذائب يترك الماء رائقًا ولا شيء في قاعه؛ وهنا الماء
عكر وقد استقرت طبقة في القاع.
\end{solution}

\begin{solution}{exo:g2:mixing-and-dissolving:5}
هي رائقة (نرى من خلالها) لكنها ليست عديمة اللون (لونها وردي). وهي
\omterm{def:g2:mixing-and-dissolving:dissolve}{محلول}: فقد يكون \omterm{def:g2:mixing-and-dissolving:dissolve}{للمحلول} لون.
\end{solution}

\begin{solution}{exo:g2:mixing-and-dissolving:6}
طعم الشاي حلو: فالسكر لا يزال فيه، منتشرًا في الماء في قطع أصغر من أن
تُرى.
\end{solution}

\begin{solution}{exo:g2:mixing-and-dissolving:7}
تجفف الشمس الماء؛ ويبقى الملح الذي كان ذائبًا في ماء
البحر.
\end{solution}

\begin{solution}{exo:g2:mixing-and-dissolving:8}
«ذاب». الانصهار هو ما يحدث لمكعب الجليد من تلقاء نفسه، دون ماء؛
أما السكر فيحتاج إلى الماء، ولا يزال موجودًا في الشاي.
\end{solution}

\begin{solution}{exo:g2:mixing-and-dissolving:9}
الكأس التي فيها الرمل عكرة وفي قاعها طبقة؛ أما ماء السكر فرائق ولا
شيء في قاعه.
\end{solution}

\begin{solution}{exo:g2:mixing-and-dissolving:10}
ماء الصنبور ليس ماءً فقط: فثمة شيء ذائب فيه. وحين يجف الماء يبقى هذا
الجسم الصلب على هيئة الحلقة البيضاء.
\end{solution}

\begin{solution}{exo:g2:mixing-and-dissolving:11}
تترك بضع قطرات من كل كأس تجف على صحن داكن نظيف. الماء النقي لا يترك
شيئًا؛ والماء المالح يترك ملحًا أبيض.
\end{solution}
